% --- Archon physics formalization source begin ---
% archon:physics
% archon:covers PhyXMiniProblems/problem_phyx_mini_0125.lean
% archon:source-report reports/phyx_mini/problem_phyx_mini_0125.source.json

\chapter{Physics problem phyx\_mini\_0125}
\label{ch:PhyXMiniProblems_problem_phyx_mini_0125}

\paragraph{Problem source.}
\#\# Physical scenario

A screen containing two slits \textbackslash{}(0.100\textbackslash{},\textbackslash{}mathrm\{mm\}\textbackslash{}) apart is \textbackslash{}(1.20\textbackslash{},\textbackslash{}mathrm\{m\}\textbackslash{}) from the viewing screen. Light of wavelength \textbackslash{}(\textbackslash{}lambda=500\textbackslash{},\textbackslash{}mathrm\{nm\}\textbackslash{}) falls on the slits from a distant source.

\#\# Question

Approximately how far apart will adjacent bright interference fringes be on the screen?

\#\# Answer choices

- A: A: 6.40mm

- B: B: 6.00mm

- C: C: 6.20mm

- D: D: 6.60mm

\#\# Figure caption (auxiliary; use the image as primary evidence)

The image illustrates the double-slit experiment schematic often used in physics to demonstrate wave interference. Here's a detailed description of the elements:

- **Slits**: Two vertical slits, labeled as \textbackslash{}( S\_1 \textbackslash{}) and \textbackslash{}( S\_2 \textbackslash{}), are shown close together on the left side. The distance between them is marked as \textbackslash{}( d \textbackslash{}).

- **Screen**: A screen is depicted on the right side, with a pattern suggesting wave interference. The screen receives wave patterns from the slits.

- **Waves**: Blue wave lines on the screen indicate the wavefronts emanating from the slits.

- **Angles**: Two angles, \textbackslash{}( \textbackslash{}theta\_1 \textbackslash{}) and \textbackslash{}( \textbackslash{}theta\_2 \textbackslash{}), are shown with dashed lines originating from \textbackslash{}( S\_1 \textbackslash{}) and \textbackslash{}( S\_2 \textbackslash{}), respectively. These represent the angles at which waves are traveling towards the screen.

- **Distances**: The horizontal distance from the slits to the screen is marked as \textbackslash{}( \textbackslash{}ell \textbackslash{}). Perpendicular distances from the central line to wave maxima on the screen are labeled \textbackslash{}( x\_1 \textbackslash{}) and \textbackslash{}( x\_2 \textbackslash{}).

- **Colors and Lines**: The wave fronts are blue, and the angles are illustrated with dashed purple lines.

Overall, the diagram is used to explain how interference patterns form on a screen as a result of waves passing through the slits and the relative phase difference between them.

\#\# Dataset metadata

Wave Optics; Physical Model Grounding Reasoning, Spatial Relation Reasoning

\paragraph{Recorded answer/context.}
C: C: 6.20mm

\paragraph{Figure/image path.}
/root/proposal\_for\_physic/science-mango/phyx\_mini\_run/phyx\_data/test\_image/125.png

\paragraph{Formalization target.}
create a compiling Lean file with sorry bodies at `PhyXMiniProblems/problem\_phyx\_mini\_0125.lean`. The Lean declarations must preserve the physical quantities, dimensions or dimensional roles, figure labels, governing-law hypotheses, and final relation expressed by this problem.
Use Mathlib/Physlib names found through LeanExplore where available. If a physics API is missing, introduce faithful local abstractions rather than scalar placeholder aliases.

\begin{theorem}[Physics formalization target]
\label{thm:physics:phyx_mini_0125:target}
\lean{PhyXMiniProblems.ProblemPhyXMini0125.problem_phyx_mini_0125}
\uses{def:physics:phyx-mini-0125:phyxminiproblems-problemphyxmini0125-doubleslitexperiment, def:physics:phyx-mini-0125:phyxminiproblems-problemphyxmini0125-hasdepictedlayout, def:physics:phyx-mini-0125:phyxminiproblems-problemphyxmini0125-hasdistantinphaseillumination, def:physics:phyx-mini-0125:phyxminiproblems-problemphyxmini0125-hasphysicalparameters, def:physics:phyx-mini-0125:phyxminiproblems-problemphyxmini0125-hasstatedreadouts, def:physics:phyx-mini-0125:phyxminiproblems-problemphyxmini0125-satisfiesconstructivebrightfringelaw, def:physics:phyx-mini-0125:phyxminiproblems-problemphyxmini0125-satisfiesstraightlinescreengeometry, def:physics:phyx-mini-0125:phyxminiproblems-problemphyxmini0125-adjacentbrightfringespacingmillimeters, def:physics:phyx-mini-0125:phyxminiproblems-problemphyxmini0125-isnearestdisplayedchoice}
The assigned autoformalize agent should translate this physics problem into Lean declarations in the covered file, with theorem and lemma proof bodies written as `by sorry`.
\end{theorem}
\begin{proof}
This is an autoformalization task, not a proof task. Produce statements that can later be proved by the physics prover without weakening the physical model.
\end{proof}

% --- Archon declaration topology begin ---
\section{Declaration topology}
\begin{definition}[Length Quantity]
  \label{def:physics:phyx-mini-0125:phyxminiproblems-problemphyxmini0125-lengthquantity}
  \lean{PhyXMiniProblems.ProblemPhyXMini0125.LengthQuantity}
  A physical length represented independently of the unit used to read it
\end{definition}
\begin{proof}
  This is the declared model definition of Length Quantity.
\end{proof}
\begin{definition}[length Value In]
  \label{def:physics:phyx-mini-0125:phyxminiproblems-problemphyxmini0125-lengthvaluein}
  \lean{PhyXMiniProblems.ProblemPhyXMini0125.lengthValueIn}
  \uses{def:physics:phyx-mini-0125:phyxminiproblems-problemphyxmini0125-lengthquantity}
  Read a physical length as a real scalar in a selected length unit
\end{definition}
\begin{proof}
  This is the declared model definition of length Value In.
\end{proof}
\begin{definition}[meters Value]
  \label{def:physics:phyx-mini-0125:phyxminiproblems-problemphyxmini0125-metersvalue}
  \lean{PhyXMiniProblems.ProblemPhyXMini0125.metersValue}
  \uses{def:physics:phyx-mini-0125:phyxminiproblems-problemphyxmini0125-lengthquantity, def:physics:phyx-mini-0125:phyxminiproblems-problemphyxmini0125-lengthvaluein}
  The metre readout of a physical length
\end{definition}
\begin{proof}
  This is the declared model definition of meters Value.
\end{proof}
\begin{definition}[millimeters Value]
  \label{def:physics:phyx-mini-0125:phyxminiproblems-problemphyxmini0125-millimetersvalue}
  \lean{PhyXMiniProblems.ProblemPhyXMini0125.millimetersValue}
  \uses{def:physics:phyx-mini-0125:phyxminiproblems-problemphyxmini0125-lengthquantity, def:physics:phyx-mini-0125:phyxminiproblems-problemphyxmini0125-lengthvaluein}
  The millimetre readout used for slit and fringe spacings
\end{definition}
\begin{proof}
  This is the declared model definition of millimeters Value.
\end{proof}
\begin{definition}[nanometers Value]
  \label{def:physics:phyx-mini-0125:phyxminiproblems-problemphyxmini0125-nanometersvalue}
  \lean{PhyXMiniProblems.ProblemPhyXMini0125.nanometersValue}
  \uses{def:physics:phyx-mini-0125:phyxminiproblems-problemphyxmini0125-lengthquantity, def:physics:phyx-mini-0125:phyxminiproblems-problemphyxmini0125-lengthvaluein}
  The nanometre readout used for the light wavelength
\end{definition}
\begin{proof}
  This is the declared model definition of nanometers Value.
\end{proof}
\begin{definition}[Slit Label]
  \label{def:physics:phyx-mini-0125:phyxminiproblems-problemphyxmini0125-slitlabel}
  \lean{PhyXMiniProblems.ProblemPhyXMini0125.SlitLabel}
  The two slit labels printed in the primary figure
\end{definition}
\begin{proof}
  This is the declared model definition of Slit Label.
\end{proof}
\begin{definition}[Bright Fringe Label]
  \label{def:physics:phyx-mini-0125:phyxminiproblems-problemphyxmini0125-brightfringelabel}
  \lean{PhyXMiniProblems.ProblemPhyXMini0125.BrightFringeLabel}
  The two upper bright maxima labeled by (θ₁, x₁) and (θ₂, x₂)
\end{definition}
\begin{proof}
  This is the declared model definition of Bright Fringe Label.
\end{proof}
\begin{definition}[order]
  \label{def:physics:phyx-mini-0125:phyxminiproblems-problemphyxmini0125-brightfringelabel-order}
  \lean{PhyXMiniProblems.ProblemPhyXMini0125.BrightFringeLabel.order}
  \uses{def:physics:phyx-mini-0125:phyxminiproblems-problemphyxmini0125-brightfringelabel}
  The constructive-interference order represented by each figure label
\end{definition}
\begin{proof}
  This is the declared model definition of order.
\end{proof}
\begin{definition}[Slit Arrangement]
  \label{def:physics:phyx-mini-0125:phyxminiproblems-problemphyxmini0125-slitarrangement}
  \lean{PhyXMiniProblems.ProblemPhyXMini0125.SlitArrangement}
  Qualitative geometry of the two openings in the slit screen
\end{definition}
\begin{proof}
  This is the declared model definition of Slit Arrangement.
\end{proof}
\begin{definition}[Screen Orientation]
  \label{def:physics:phyx-mini-0125:phyxminiproblems-problemphyxmini0125-screenorientation}
  \lean{PhyXMiniProblems.ProblemPhyXMini0125.ScreenOrientation}
  Qualitative orientation of the viewing screen relative to the optical axis
\end{definition}
\begin{proof}
  This is the declared model definition of Screen Orientation.
\end{proof}
\begin{definition}[Incident Wavefront Regime]
  \label{def:physics:phyx-mini-0125:phyxminiproblems-problemphyxmini0125-incidentwavefrontregime}
  \lean{PhyXMiniProblems.ProblemPhyXMini0125.IncidentWavefrontRegime}
  The wavefront regime of the light incident on the two slits
\end{definition}
\begin{proof}
  This is the declared model definition of Incident Wavefront Regime.
\end{proof}
\begin{definition}[Double Slit Experiment]
  \label{def:physics:phyx-mini-0125:phyxminiproblems-problemphyxmini0125-doubleslitexperiment}
  \lean{PhyXMiniProblems.ProblemPhyXMini0125.DoubleSlitExperiment}
  \uses{def:physics:phyx-mini-0125:phyxminiproblems-problemphyxmini0125-lengthquantity, def:physics:phyx-mini-0125:phyxminiproblems-problemphyxmini0125-slitlabel, def:physics:phyx-mini-0125:phyxminiproblems-problemphyxmini0125-brightfringelabel, def:physics:phyx-mini-0125:phyxminiproblems-problemphyxmini0125-slitarrangement, def:physics:phyx-mini-0125:phyxminiproblems-problemphyxmini0125-screenorientation, def:physics:phyx-mini-0125:phyxminiproblems-problemphyxmini0125-incidentwavefrontregime}
  The physical quantities and figure-derived data for the experiment. Phase remains abstract because the source gives no phase unit or scalar calibration. Only equality of the incident phases at S₁ and S₂ is used. The fields fringeAngleRadians and fringeOffsetFromCentral interpret the figure pairs (θ₁, x₁) and (θ₂, x₂)
\end{definition}
\begin{proof}
  This is the declared model definition of Double Slit Experiment.
\end{proof}
\begin{definition}[Has Depicted Layout]
  \label{def:physics:phyx-mini-0125:phyxminiproblems-problemphyxmini0125-hasdepictedlayout}
  \lean{PhyXMiniProblems.ProblemPhyXMini0125.HasDepictedLayout}
  \uses{def:physics:phyx-mini-0125:phyxminiproblems-problemphyxmini0125-doubleslitexperiment}
  The categorical slit and screen layout visible in the primary figure
\end{definition}
\begin{proof}
  This is the declared model definition of Has Depicted Layout.
\end{proof}
\begin{definition}[Has Distant In Phase Illumination]
  \label{def:physics:phyx-mini-0125:phyxminiproblems-problemphyxmini0125-hasdistantinphaseillumination}
  \lean{PhyXMiniProblems.ProblemPhyXMini0125.HasDistantInPhaseIllumination}
  \uses{def:physics:phyx-mini-0125:phyxminiproblems-problemphyxmini0125-doubleslitexperiment}
  The distant monochromatic source is modeled by a planar incident wavefront and equal incident phase at the two slits. This does not constrain any fringe position on the viewing screen
\end{definition}
\begin{proof}
  This is the declared model definition of Has Distant In Phase Illumination.
\end{proof}
\begin{definition}[Has Physical Parameters]
  \label{def:physics:phyx-mini-0125:phyxminiproblems-problemphyxmini0125-hasphysicalparameters}
  \lean{PhyXMiniProblems.ProblemPhyXMini0125.HasPhysicalParameters}
  \uses{def:physics:phyx-mini-0125:phyxminiproblems-problemphyxmini0125-metersvalue, def:physics:phyx-mini-0125:phyxminiproblems-problemphyxmini0125-brightfringelabel, def:physics:phyx-mini-0125:phyxminiproblems-problemphyxmini0125-doubleslitexperiment}
  Positivity and angular-range conditions for the depicted experiment
\end{definition}
\begin{proof}
  This is the declared model definition of Has Physical Parameters.
\end{proof}
\begin{definition}[Has Stated Readouts]
  \label{def:physics:phyx-mini-0125:phyxminiproblems-problemphyxmini0125-hasstatedreadouts}
  \lean{PhyXMiniProblems.ProblemPhyXMini0125.HasStatedReadouts}
  \uses{def:physics:phyx-mini-0125:phyxminiproblems-problemphyxmini0125-metersvalue, def:physics:phyx-mini-0125:phyxminiproblems-problemphyxmini0125-millimetersvalue, def:physics:phyx-mini-0125:phyxminiproblems-problemphyxmini0125-nanometersvalue, def:physics:phyx-mini-0125:phyxminiproblems-problemphyxmini0125-doubleslitexperiment}
  Numerical readouts stated in the problem: d = 0.100 mm, ℓ = 1.20 m, and λ = 500 nm. No angle or fringe-spacing answer occurs here
\end{definition}
\begin{proof}
  This is the declared model definition of Has Stated Readouts.
\end{proof}
\begin{definition}[Satisfies Constructive Bright Fringe Law]
  \label{def:physics:phyx-mini-0125:phyxminiproblems-problemphyxmini0125-satisfiesconstructivebrightfringelaw}
  \lean{PhyXMiniProblems.ProblemPhyXMini0125.SatisfiesConstructiveBrightFringeLaw}
  \uses{def:physics:phyx-mini-0125:phyxminiproblems-problemphyxmini0125-metersvalue, def:physics:phyx-mini-0125:phyxminiproblems-problemphyxmini0125-brightfringelabel, def:physics:phyx-mini-0125:phyxminiproblems-problemphyxmini0125-doubleslitexperiment}
  Constructive double-slit interference law for each labeled bright fringe: d sin θₘ = m λ. The equation is written using metre readouts, so all length-valued factors use the same unit
\end{definition}
\begin{proof}
  This is the declared model definition of Satisfies Constructive Bright Fringe Law.
\end{proof}
\begin{definition}[Satisfies Straight Line Screen Geometry]
  \label{def:physics:phyx-mini-0125:phyxminiproblems-problemphyxmini0125-satisfiesstraightlinescreengeometry}
  \lean{PhyXMiniProblems.ProblemPhyXMini0125.SatisfiesStraightLineScreenGeometry}
  \uses{def:physics:phyx-mini-0125:phyxminiproblems-problemphyxmini0125-metersvalue, def:physics:phyx-mini-0125:phyxminiproblems-problemphyxmini0125-brightfringelabel, def:physics:phyx-mini-0125:phyxminiproblems-problemphyxmini0125-doubleslitexperiment}
  Exact straight-line screen geometry for each labeled maximum: xₘ = ℓ tan θₘ. This generic law contains no requested numerical spacing
\end{definition}
\begin{proof}
  This is the declared model definition of Satisfies Straight Line Screen Geometry.
\end{proof}
\begin{definition}[adjacent Bright Fringe Spacing Millimeters]
  \label{def:physics:phyx-mini-0125:phyxminiproblems-problemphyxmini0125-adjacentbrightfringespacingmillimeters}
  \lean{PhyXMiniProblems.ProblemPhyXMini0125.adjacentBrightFringeSpacingMillimeters}
  \uses{def:physics:phyx-mini-0125:phyxminiproblems-problemphyxmini0125-millimetersvalue, def:physics:phyx-mini-0125:phyxminiproblems-problemphyxmini0125-doubleslitexperiment}
  The measured separation x₂ - x₁ of the adjacent labeled bright fringes
\end{definition}
\begin{proof}
  This is the declared model definition of adjacent Bright Fringe Spacing Millimeters.
\end{proof}
\begin{definition}[Answer Choice]
  \label{def:physics:phyx-mini-0125:phyxminiproblems-problemphyxmini0125-answerchoice}
  \lean{PhyXMiniProblems.ProblemPhyXMini0125.AnswerChoice}
  Labels of the four displayed answer choices
\end{definition}
\begin{proof}
  This is the declared model definition of Answer Choice.
\end{proof}
\begin{definition}[answer Spacing Millimeters]
  \label{def:physics:phyx-mini-0125:phyxminiproblems-problemphyxmini0125-answerspacingmillimeters}
  \lean{PhyXMiniProblems.ProblemPhyXMini0125.answerSpacingMillimeters}
  \uses{def:physics:phyx-mini-0125:phyxminiproblems-problemphyxmini0125-answerchoice}
  The fringe spacing printed beside each answer choice, in millimetres
\end{definition}
\begin{proof}
  This is the declared model definition of answer Spacing Millimeters.
\end{proof}
\begin{definition}[recorded Answer Choice]
  \label{def:physics:phyx-mini-0125:phyxminiproblems-problemphyxmini0125-recordedanswerchoice}
  \lean{PhyXMiniProblems.ProblemPhyXMini0125.recordedAnswerChoice}
  \uses{def:physics:phyx-mini-0125:phyxminiproblems-problemphyxmini0125-answerchoice}
  Dataset metadata only: the source records answer choice C
\end{definition}
\begin{proof}
  This is the declared model definition of recorded Answer Choice.
\end{proof}
\begin{definition}[Is Nearest Displayed Choice]
  \label{def:physics:phyx-mini-0125:phyxminiproblems-problemphyxmini0125-isnearestdisplayedchoice}
  \lean{PhyXMiniProblems.ProblemPhyXMini0125.IsNearestDisplayedChoice}
  \uses{def:physics:phyx-mini-0125:phyxminiproblems-problemphyxmini0125-doubleslitexperiment, def:physics:phyx-mini-0125:phyxminiproblems-problemphyxmini0125-adjacentbrightfringespacingmillimeters, def:physics:phyx-mini-0125:phyxminiproblems-problemphyxmini0125-answerchoice, def:physics:phyx-mini-0125:phyxminiproblems-problemphyxmini0125-answerspacingmillimeters}
  A choice is correct when its displayed value is nearest to the physical spacing
\end{definition}
\begin{proof}
  This is the declared model definition of Is Nearest Displayed Choice.
\end{proof}
% --- Archon declaration topology end ---
% --- Archon physics formalization source end ---
